\documentclass[11pt]{article}
\usepackage{paper-style}

\title{\textbf{Finite-Complement Congruence Topologies for Partial Algebras}}
\author{Adrian Puha}
\date{21 August 2026}

\begin{document}
\maketitle

\begin{abstract}
Let $D$ be a partial algebra with carrier $A$, and let $F(D)$ be its free
compatible total extension.  We consider compatible total extensions that fix
$A$ pointwise and contain only finitely many elements outside $A$.  Their
kernels define a descending family of congruences on $F(D)$ and hence a
Hausdorff congruence topology.

A preliminary observation is that every finite subset of $F(D)$ is
discriminated by one such extension.  The proof is a Rees-congruence
construction: retain the finite subterm closure and collapse its complement to
a single class.  This is the algebraic form of the standard finite-subterm
plus sink construction from tree automata.

When $A$ is finite, the resulting topology is the ordinary profinite topology.
The case of interest is therefore an infinite base, where the allowed target
algebras may themselves be infinite.  We give a sufficient condition for the
Hausdorff completion to be strictly larger than $F(D)$.  If $D$ has an
undefined base application and admits a nontrivial unary context acting as the
identity on every element of $A$, then iterating that context from the
undefined term yields a Cauchy sequence with no limit in $F(D)$.  Natural
arithmetic with division undefined at zero gives a simple example.
\end{abstract}

\section{Introduction}\label{sec:introduction}

A partial algebra is an algebra in which some applications of the basic
operations are undefined.  A classical way to compare a partial algebra with
total algebras is to embed it into, or map it to, compatible total algebras.
We study a restricted version of this problem in which the original carrier is
kept pointwise and only finitely many new elements are permitted.

Let $D$ be a one-sorted partial algebra with carrier $A$ and a finite binary
signature $\Sigma$.  We use the standard free compatible completion $\Free D$
of $D$: applications already defined in $D$ are reduced to their prescribed
values, while undefined applications remain as finite formal terms.  The free
completion is background; the question considered here is how finite subsets
of $\Free D$ can be distinguished by compatible total completions of $D$.

\begin{definition}
A \emph{finite-complement completion} of $D$ is a total $\Sigma$-algebra $T$
containing $A$ as a set, such that
\begin{enumerate}
  \item every operation value already defined in $D$ has the same value in
        $T$, and
  \item $T\setminus A$ is finite.
\end{enumerate}
The universal property of $\Free D$ gives a canonical homomorphism
$\pi_T:\Free D\to T$ fixing $A$ pointwise.
\end{definition}

The main result says that the family of all such completions discriminates every
finite subset of the free completion.

\begin{theorem}[Finite-subset discrimination; informal form]
\label{thm:intro-discrimination}
For every finite subset $X\subseteq\Free D$ there exists a finite-complement
completion $T$ of $D$ such that $\pi_T$ is injective on $X$.
\end{theorem}

The proof is elementary and constructive.  Take all subterms occurring in the
chosen finite set.  Keep base elements unchanged, assign one new element to
each non-base subterm that must be remembered, and send all remaining cases to
a single sink.  The normal-form property of the free completion guarantees
that this finite operation table extends the partial operations of $D$.

This statement resembles the standard passage from residual separation to
finite-set discrimination.  There is, however, a simple obstruction to the
usual proof by finite products.  If $A$ is infinite, the class of
finite-complement completions need not be closed under products over the fixed
copy of $A$.  For example, each inclusion
\[
 A\hookrightarrow A\sqcup\{\ast\}
\]
has one new element, whereas the diagonal copy of $A$ in
$(A\sqcup\{\ast\})^2$ has infinitely many elements outside it whenever $A$ is
infinite.  Thus the direct finite-subterm construction is needed if one wants
to remain inside the same class of completions.

We then use these completions to define finite-complement congruences on
$\Free D$.  For $m\ge 1$, write $x\fceq{m}y$ when every finite-complement
completion with at most $m$ new elements identifies $x$ and $y$.  The
finite-subset theorem implies
\[
 \bigcap_{m\ge1}\fceq{m}\;=\;\Delta_{\Free D}.
\]
Hence these congruences define a separated uniform structure, and we consider
its Cauchy completion.

The second part of the paper gives two criteria ensuring that this completion
contains new points.  For finite $A$, properness is equivalent to genuine
partiality of $D$.  For arbitrary $A$, we give a sufficient condition using a
unary context that fixes every base element.  The proof is again elementary:
finite-state periodicity produces a Cauchy subsequence, while a finite-subterm
separator rules out convergence to any element of $\Free D$.

The scope of the contribution is deliberately narrow.  Free completions of
partial algebras, residuality, discrimination, finite-state periodicity and
Cauchy completion are established ideas.  The point studied here is their
interaction under the additional requirement that a possibly infinite base be
retained pointwise while only its complement is finite.  The literature
comparison in \cref{sec:related-work} is written accordingly.

All proofs are given in ordinary mathematical language.  A machine-checked
formalization is mentioned only in an appendix as supplementary verification
and is not used as a premise in any argument.

\section{Free compatible completion}\label{sec:free-completion}

We recall the concrete free completion used throughout.  Nothing in this
section is claimed as new.

Let $D=(A,\Sigma,\delta)$ be a one-sorted partial algebra with binary signature
$\Sigma$.  For $f\in\Sigma$ and $a,b\in A$, the value $\delta_f(a,b)$ is either
an element of $A$ or is undefined.

A \emph{normal term} over $D$ is defined recursively.  Every $a\in A$ is a
normal term.  If $s,t$ are normal terms and $f\in\Sigma$, then $f(s,t)$ is a
normal term unless both $s,t$ lie in $A$ and $\delta_f(s,t)$ is defined; in
that exceptional case $f(s,t)$ is replaced by the corresponding base value.
Let $\Free D$ denote the set of normal terms.  Each basic operation is made
total on $\Free D$ by forming the formal application and then applying this
single normalization rule.

The inclusion $\eta:A\hookrightarrow\Free D$ is injective, and every operation
value already defined in $D$ is preserved:
\[
 \delta_f(a,b)=c
 \quad\Longrightarrow\quad
 f_{\Free D}(a,b)=c.
\]

A total $\Sigma$-algebra $T$ together with a map $i:A\to T$ is
\emph{compatible with $D$} when
\[
 \delta_f(a,b)=c
 \quad\Longrightarrow\quad
 f_T(i(a),i(b))=i(c).
\]
Evaluation of normal terms in $T$ gives the standard universal property.

\begin{theorem}[Free compatible completion]\label{thm:free-completion}
For every compatible pair $(T,i)$ there is a unique homomorphism
$\pi_T:\Free D\to T$ such that $\pi_T\eta=i$.
\end{theorem}

\begin{proof}
Define $\pi_T$ recursively.  On $A$ put $\pi_T(a)=i(a)$.  For a non-base normal
term $f(s,t)$ put
\[
 \pi_T(f(s,t))=f_T(\pi_T(s),\pi_T(t)).
\]
The only possible ambiguity would come from an application of $f$ to two base
values on which $D$ is already defined.  Such an application is not a
non-base normal term: by definition it has already been reduced to its base
value.  Compatibility therefore makes the recursive definition agree with all
reductions carried out in $\Free D$.  The homomorphism property is immediate,
and uniqueness follows by induction on term formation.
\end{proof}

We shall repeatedly use the following elementary consequence of normality.

\begin{lemma}[Normality]\label{lem:normality}
If $f(s,t)\in\Free D\setminus A$, then it is not the case that
$s,t\in A$ and $\delta_f(s,t)$ is defined.
\end{lemma}

This is exactly the local fact needed in the finite-complement construction of
the next section.

\section{Finite-complement separation}\label{sec:finite-separation}

For a normal term $x\in F(D)$, let $\operatorname{Sub}(x)$ denote its finite
set of normal subterms.  If $X\subseteq F(D)$ is finite, set
\[
 S_X=\bigcup_{x\in X}\operatorname{Sub}(x),
 \qquad
 N_X=S_X\setminus A.
\]
Thus $N_X$ is finite and $S_X$ is closed under taking subterms.

\begin{lemma}[Finite-complement Rees quotient]
\label{lem:finite-complement-rees}
For every finite $X\subseteq F(D)$ there is a congruence $\rho_X$ on $F(D)$
such that every element of $A\cup N_X$ is a singleton $\rho_X$-class and, if
nonempty, $F(D)\setminus(A\cup N_X)$ is one additional class.  Consequently
$F(D)/\rho_X$ is a finite-complement compatible extension of $D$ and the
quotient map is injective on $X$.
\end{lemma}

\begin{proof}
Put
\[
 I_X=F(D)\setminus(A\cup N_X).
\]
We first observe that $I_X$ is absorbing in every coordinate: if one argument
of a basic operation lies in $I_X$, then so does the result.  Indeed, let
$f(s,t)$ be the result of applying a basic operation and suppose, for example,
that $s\in I_X$.  The result cannot lie in $A$, because reduction to a base
value is possible only when both arguments lie in $A$.  Nor can it lie in
$N_X$: if the normal term $f(s,t)$ were a selected subterm, then its subterm
$s$ would belong to $S_X$, contrary to $s\in I_X$.  The same argument applies
when $t\in I_X$.

Define
\[
 \rho_X=\Delta_{F(D)}\cup(I_X\times I_X).
\]
The absorption property makes $\rho_X$ a congruence; when $I_X$ is nonempty,
this is the Rees congruence determined by the absorbing subalgebra $I_X$
\cite{Tichy1981}.  Every base element remains a singleton class.  Outside the
image of $A$ the quotient has the singleton classes indexed by $N_X$, together
with at most one additional class represented by $I_X$.  Hence
\[
 |(F(D)/\rho_X)\setminus A|\le |N_X|+1<\infty.
\]
Since $X\subseteq A\cup N_X$, the quotient map is injective on $X$.
\end{proof}

The construction is the algebraic form of the familiar finite-subterm
construction with one sink state in tree automata \cite{ComonEtAl2008}.

\begin{corollary}[Separation]
\label{cor:finite-complement-separation}
If $x\ne y$ in $F(D)$, then there is a finite-complement compatible extension
$T$ for which the canonical map $F(D)\to T$ separates $x$ and $y$.
\end{corollary}

\begin{remark}
When $A$ is infinite, finite-complement extensions are not in general closed
under ordinary products over the diagonal copy of $A$: already the diagonal
copy of $A$ in $(A\sqcup\{\ast\})^2$ has infinitely many points outside it.
This observation is not needed in the proof above; it merely explains why the
usual product-of-separators argument does not stay inside the same class.
\end{remark}

\section{The induced congruence topology}\label{sec:completion}

For $m\ge1$, define a congruence on $F(D)$ by
\[
 \theta_m=
 \bigcap\{\ker(\pi_T):
 T\text{ is a finite-complement compatible extension and }
 |T\setminus A|\le m\}.
\]
Equivalently, $x\equiv y\pmod{\theta_m}$ when every such extension identifies
$x$ and $y$.  Each $\theta_m$ is a congruence and
\[
 \theta_{m+1}\subseteq\theta_m.
\]
By \cref{cor:finite-complement-separation},
\[
 \bigcap_{m\ge1}\theta_m=\Delta_{F(D)}.
\]
Hence the cosets of the $\theta_m$ form a countable Hausdorff congruence
topology on $F(D)$; this is a standard construction for universal algebras
\cite{BulmanFleming1971}.  We denote its Hausdorff completion by
$\widehat{F(D)}_{\mathrm{fc}}$.

\subsection{Finite bases}

When $A$ is finite, the preceding topology is not a new kind of completion.

\begin{proposition}[Finite base gives the profinite topology]
\label{prop:finite-base-profinite}
Assume that $A$ and $\Sigma$ are finite.  Then the topology generated by the
congruences $\theta_m$ is the ordinary profinite topology on $F(D)$.
\end{proposition}

\begin{proof}
First, each $\theta_m$ has finite index.  Up to relabelling of the finitely many
points outside the fixed finite set $A$, there are only finitely many total
$\Sigma$-algebra structures on carriers of size at most $|A|+m$.  Thus only
finitely many kernels occur in the defining intersection for $\theta_m$, and
an intersection of finitely many finite-index congruences again has finite
index.  Hence every $\theta_m$ is an entourage for the profinite topology.

Conversely, let $\alpha$ be any finite-index congruence on $F(D)$.  Apply
\cref{lem:finite-complement-rees} to the finite set $X=A$.  This gives a
finite-index congruence $\beta$ whose restriction to $A$ is equality.  Put
\[
 \gamma=\alpha\cap\beta.
\]
Then $\gamma$ has finite index, $\gamma\subseteq\alpha$, and $A$ remains
pointwise distinct in $F(D)/\gamma$.  After identifying the image of $A$ with
$A$ itself, the quotient $F(D)/\gamma$ is a finite-complement compatible
extension.  If its complement has size $m$, then
\[
 \theta_m\subseteq\gamma\subseteq\alpha.
\]
Thus the family $(\theta_m)$ is cofinal among the finite-index congruences, so
it induces exactly the profinite topology.
\end{proof}

The remainder of the section concerns the situation in which $A$ may be
infinite.  Then an allowed target $T$ can be infinite, and the congruences
$\theta_m$ need not have finite index.

\subsection{A criterion over an arbitrary base}

We use one elementary observation about finite deterministic orbits.

\begin{lemma}\label{lem:finite-orbit-factorials}
Let $s_0,s_1,\ldots$ be an orbit of a self-map.  If the orbit takes at most
$N$ distinct values, then $s_{k!}$ is independent of $k$ for all sufficiently
large $k$.  The threshold can be chosen using only $N$.
\end{lemma}

\begin{proof}
Among the first $N+1$ positions two coincide.  The orbit is therefore eventually
periodic, with preperiod and period bounded by $N$; every sufficiently large
factorial lies beyond the preperiod and is divisible by the period.
\end{proof}

A \emph{one-hole unary context} $C[-]$ is a term with one distinguished
occurrence of a variable.  We say that it \emph{fixes the base} if substitution
of any $a\in A$ is defined in $D$ and evaluates to $a$:
\[
 C^D(a)=a\qquad(a\in A).
\]
We call the context nontrivial when it contains at least one operation symbol.

\begin{theorem}[Infinite-base completion criterion]
\label{thm:infinite-base-criterion}
Suppose that some basic application in $D$ is undefined, and let $c_0\in
F(D)\setminus A$ be the corresponding normal term.  If there is a nontrivial
one-hole unary context $C[-]$ that fixes the base, then the canonical map
\[
 F(D)\longrightarrow\widehat{F(D)}_{\mathrm{fc}}
\]
is not surjective.
\end{theorem}

\begin{proof}
Define
\[
 c_{n+1}=C[c_n]\qquad(n\ge0).
\]
Because $C[-]$ is nontrivial and $c_0$ is non-base, no normalization step can
remove the occurrence of $c_n$ inside $C[c_n]$; hence the construction depth
strictly increases and the terms $c_n$ are pairwise distinct.

Consider the subsequence $d_k=c_{k!}$.  Fix $m$ and a finite-complement
compatible extension $T$ with $|T\setminus A|\le m$.  The images of the $c_n$
form an orbit under the unary operation induced by $C[-]$ on $T$.  If this
orbit enters $A$, it is constant from then on because $C$ fixes every element
of $A$.  Otherwise it remains in the at most $m$ elements of $T\setminus A$.
Thus in either case the orbit has at most $m+1$ values before becoming
periodic.  By \cref{lem:finite-orbit-factorials}, for all sufficiently large
$k,l$, with a threshold depending only on $m$,
\[
 \pi_T(d_k)=\pi_T(d_l).
\]
Since the threshold is uniform over all such $T$, we have
$d_k\equiv d_l\pmod{\theta_m}$ for all sufficiently large $k,l$.  As $m$ was
arbitrary, $(d_k)$ is Cauchy.

It remains to show that this Cauchy sequence has no limit in $F(D)$.  Fix
$x\in F(D)$.  Choose $N$ so large that the construction depth of $c_N$ exceeds
that of $x$, and apply \cref{lem:finite-complement-rees} to
$X=\{x,c_N\}$.  In the resulting quotient, $x$ and every selected subterm are
singleton classes.  The next term $c_{N+1}=C[c_N]$ is not among those selected
subterms: it has greater depth than both selected roots.  Hence it belongs to
the collapsed ideal class.  Since that class is absorbing under every basic
operation, all later $c_j$ also belong to it.  Therefore this one
finite-complement quotient separates $x$ from every sufficiently late $d_k$.
So $(d_k)$ cannot converge to $x$.

Since $x$ was arbitrary, the Cauchy sequence $(d_k)$ has no limit in $F(D)$.
Its limit in the Hausdorff completion is therefore a point outside the image of
$F(D)$.
\end{proof}

\begin{remark}
The theorem is only a sufficient criterion.  It does not characterize when the
finite-complement congruence topology over an infinite base is incomplete.
The point of the hypothesis $C^D|_A=\mathrm{id}_A$ is that, although the target
algebra may be infinite, the particular orbit used in the proof cannot employ
the retained base as unbounded memory.
\end{remark}

\section{Examples}\label{sec:examples}

The preceding criterion has elementary infinite-base examples.  They are
included only to illustrate the construction; no claim about a new arithmetic
system is intended.

\begin{example}[Natural numbers]
Let $D$ have carrier $\mathbb N$ with the usual addition and multiplication and
with division undefined when the denominator is zero.  In the free completion,
$a/0$ is therefore a non-base normal term rather than a natural number.
Addition satisfies
\[
 n+0=n\qquad(n\in\mathbb N).
\]
Taking the undefined term $0/0$ as $c_0$ and the context $x\mapsto x+0$ in
\cref{thm:base-fixing-properness} shows that
$\FCcomp{\Free D}$ is proper.
\end{example}

\begin{example}[Integers]
The same argument applies to $\mathbb Z$ with the usual ring operations and
partial division.  Again $x\mapsto x+0$ fixes the whole base, so the
finite-complement uniform completion is proper.
\end{example}

These examples should not be confused with totalized fields, meadows, common
meadows, or wheels, which impose algebraic conventions for singular division
\cite{BergstraHirshfeldTucker2009,BergstraPonseCommonMeadows,Carlstrom2004}.
Here no numerical value is assigned to $a/0$; the undefined application is
retained as a formal element of the free completion.

\section{Conclusion}

We studied a restricted class of total completions of a partial algebra: the
original carrier is preserved pointwise and only finitely many new elements may
be added.  The main theorem shows that this class discriminates every finite
subset of the free compatible completion by a single homomorphism.  The proof
is direct and uses only the finite set of subterms occurring in the elements to
be distinguished.

When the base is infinite, this statement is not obtained from pairwise
separation by the usual finite-product argument, because finite complement over
the diagonal copy of the base is not preserved by products.  This explains why
the direct finite-subterm construction is the relevant one in the present
setting.

The same family of finite-complement completions defines a separated uniformity
on the free completion.  For finite bases its Cauchy completion is proper
exactly when the original algebra is genuinely partial.  For arbitrary bases,
a context fixing every base element gives a simple sufficient condition for
properness.

The intended contribution is modest: it is the finite-complement restriction
and the resulting discrimination and completion statements.  The surrounding
free-completion, residual, finite-state and Cauchy machinery is standard and has
been presented as such throughout.


\begingroup
\small
\setlength{\bibsep}{0.25em}
\bibliographystyle{abbrvnat}
\bibliography{../references,../references-completion,../references-separation,../references-standardization}
\endgroup

\end{document}
